\documentclass[11pt]{article}

\usepackage[a4paper,margin=1in]{geometry}
\usepackage[T1]{fontenc}
\usepackage[utf8]{inputenc}
\usepackage{amsmath,amssymb,amsthm}
\usepackage{booktabs}
\usepackage{hyperref}
\usepackage{xcolor}
\usepackage{listings}
\usepackage{microtype}

\hypersetup{
  colorlinks=true,
  linkcolor=blue!50!black,
  urlcolor=blue!50!black,
  citecolor=blue!50!black
}

\lstset{
  basicstyle=\ttfamily\small,
  columns=fullflexible,
  keepspaces=true,
  frame=single,
  framerule=0.3pt,
  breaklines=true,
  showstringspaces=false
}

\newcommand{\sandhi}{\textit{sandhi}}
\newcommand{\vichcheda}{\textit{vichcheda}}
\newtheorem{proposition}{Proposition}
\newtheorem{theorem}{Theorem}
\newtheorem{lemma}{Lemma}
\newtheorem{conjecture}{Conjecture}
\newtheorem{remark}{Remark}

\title{Recursive Sandhi and Vichcheda Canonicalization\\
for Exterior-Product Contractions}
\author{Akash Dwivedi\\
\texttt{akash.d.dwivedi@gmail.com}\\
\url{https://www.linkedin.com/in/akash-dvd-3259402a/}}
\date{\today}

\begin{document}

\maketitle

\begin{abstract}
Symbolic exterior-algebra calculations must preserve both antisymmetry and the
scalar factors created when left-contraction panels share geometric factors.
This paper studies recursive stitching for contraction panels whose down
arguments are decomposable blades. The local grade-\(r\) seam identity is
proved by an iterated-contraction minor expansion: a grade-\(r\) contractor
can be stitched across a grade-\(r\) seam with scalar coefficient given by
their Capelli pairing. Nested stitching accumulates contractor and seam blades
together, while a grade-balance gate determines whether to descend, emit a
scalar, or vanish under the sole-seam invariant.

The local identity and the fixed-schedule induction theorem are proved for the
well-formed panel grammar using a level-local visibility closure lemma.
Schedule-independent normalization remains open; cases through nesting depth
six have been checked by direct symbolic expansion.
\end{abstract}

\noindent\textbf{Keywords:}
geometric algebra; exterior algebra; left contraction; symbolic
canonicalization; \sandhi; \vichcheda; Capelli; normal form

\section*{Introduction}

Exterior products provide a compact representation of antisymmetric
multi-vectors, but symbolic expressions become difficult to manipulate when
left-contraction panels are wedged together and share factors. The elementary
pattern is
\[
(a\mathbin{\lrcorner}(b\wedge s))\wedge
 (a\mathbin{\lrcorner}(s\wedge c)).
\]
The shared factor \(s\) is the seam. The base identity combines the panels as
\[
(a\cdot s)\,
 a\mathbin{\lrcorner}(b\wedge s\wedge c).
\]
The sign is determined by orienting the seam toward the join.

This paper develops the grade-matched generalization in which a grade-\(r\)
contractor meets a grade-\(r\) seam blade. The local identity is proved by an
iterated-contraction minor expansion, and nested \sandhi\ is formulated as
repeated application of this local theorem. The contractor and seam grow
together so that their pairing remains scalar.

The local grade-\(r\) identity and the fixed-schedule induction theorem are
proved for the well-formed panel grammar using a level-local visibility closure
lemma. Schedule-independent normalization remains open; cases through
nesting depth six have been checked by direct symbolic expansion.

\section{Setting}

Let \(V\) be a vector space with a nondegenerate bilinear form. We use
\(\wedge\) for the exterior product, \(\mathbin{\lrcorner}\) for left
contraction, and \(\mid\) for the scalar pairing of equal-grade blades.

For a vector \(u\),
\[
u\mathbin{\lrcorner}(X\wedge Y)
 =(u\mathbin{\lrcorner}X)\wedge Y
 +(-1)^{\operatorname{grade}(X)}
   X\wedge(u\mathbin{\lrcorner}Y).
\]

For a decomposable blade
\[
A=a_1\wedge\cdots\wedge a_r,
\]
left contraction is the ordered iteration
\[
A\mathbin{\lrcorner}X
 =
a_1\mathbin{\lrcorner}\bigl(
a_2\mathbin{\lrcorner}(\cdots
(a_r\mathbin{\lrcorner}X)\cdots)\bigr).
\]
Equivalently,
\[
(u\wedge A)\mathbin{\lrcorner}X
 =u\mathbin{\lrcorner}(A\mathbin{\lrcorner}X).
\]
This recursive rule fixes the contraction-order sign convention.
The input fragment consists of exterior products of contraction panels. Their
ordering is part of the exterior-algebra data: reordering factors contributes
parity signs, and repeated factors can make an expression zero.

\section{Base Identity}

Let \(a\) be a vector and let \(A,B\) be decomposable blades with one
distinguished shared one-vector \(s\). After reordering their factors, write
\[
  A=(-1)^{\sigma_A}U\wedge s,\qquad
  B=(-1)^{\sigma_B}s\wedge V,
  \qquad \sigma_A,\sigma_B\in\{0,1\},
\]
where \(U\) and \(V\) have no common factors. Define the oriented concatenation
\[
  A\sqcup B :=
  (-1)^{\sigma_A+\sigma_B}U\wedge s\wedge V .
\]

\begin{proposition}[Sole-shared-seam identity]
\label{prop:base-identity}
Under these hypotheses,
\[
  (a\mathbin{\lrcorner} A)\wedge(a\mathbin{\lrcorner} B)
  =
  (a\cdot s)\,
  \bigl(a\mathbin{\lrcorner}(A\sqcup B)\bigr).
\]
\end{proposition}

\begin{proof}
Put \(p=\operatorname{grade}(U)\) and
\(\epsilon=(-1)^{\sigma_A+\sigma_B}\). The graded Leibniz rule gives
\[
\begin{aligned}
a\mathbin{\lrcorner}A
 &=(-1)^{\sigma_A}
   \bigl((a\mathbin{\lrcorner}U)\wedge s
   +(-1)^p(a\cdot s)U\bigr),\\
a\mathbin{\lrcorner}B
 &=(-1)^{\sigma_B}
   \bigl((a\cdot s)V-s\wedge(a\mathbin{\lrcorner}V)\bigr).
\end{aligned}
\]
Wedging and using \(s\wedge s=0\) yields
\[
\begin{aligned}
(a\mathbin{\lrcorner}A)\wedge(a\mathbin{\lrcorner}B)
={}&\epsilon(a\cdot s)\bigl[
  (a\mathbin{\lrcorner}U)\wedge s\wedge V
 +(-1)^p(a\cdot s)U\wedge V\\
&\hspace{2.2cm}
 -(-1)^pU\wedge s\wedge(a\mathbin{\lrcorner}V)
\bigr].
\end{aligned}
\]
The bracket is exactly the following Leibniz expansion:
\[
\begin{aligned}
a\mathbin{\lrcorner}(U\wedge s\wedge V)
={}&(a\mathbin{\lrcorner}U)\wedge s\wedge V\\
&+(-1)^pU\wedge
\bigl((a\cdot s)V-s\wedge(a\mathbin{\lrcorner}V)\bigr).
\end{aligned}
\]
Since \(A\sqcup B=\epsilon(U\wedge s\wedge V)\), the claimed identity follows.
\end{proof}

\begin{remark}
The disjointness of \(U\) and \(V\) identifies \(s\) as the sole seam and
ensures a nondegenerate concatenation. If \(U\) and \(V\) share a factor,
then \(U\wedge s\wedge V=0\); the same expansion above shows that the left
side also vanishes. Thus disjointness is needed for the sole-seam
interpretation, not for the formal equality after repeated factors are
allowed.
\end{remark}

\section{Grade-\(r\) Theorem}

The scalar in the recursive construction is determined by grade. Let \(A\)
and \(S\) be decomposable blades of the same grade \(r\), and write
\[
  B=\epsilon_B\,b\wedge S,\qquad
  C=\epsilon_C\,S\wedge c,
  \qquad \epsilon_B,\epsilon_C\in\{+1,-1\},
\]
where the factors of \(b\) and \(c\) are disjoint from one another and from
the displayed seam factors. Define
\[
  B\sqcup C :=
  \epsilon_B\epsilon_C\,b\wedge S\wedge c .
\]

\begin{lemma}[Canonical orientation]
\label{lem:canonical-orientation}
Fix any total order on atom labels, and extend it recursively (for example,
lexicographically) to the immediate factors at the current panel node. For a
panel with sole seam \(S\), sorting its factors
with the seam factors adjacent to the join gives a unique aligned form:
\[
 B=\epsilon_B\,b\wedge S,\qquad
 C=\epsilon_C\,S\wedge c,
\]
where the displayed factors are in the fixed induced order. The coefficients
\(\epsilon_B,\epsilon_C\in\{+1,-1\}\) are the parities of the unique
permutations from the original factor lists to these aligned lists.
\end{lemma}

\begin{proof}
The total order fixes the order of the residual factors and of the seam
factors. The sole-seam hypothesis fixes which factors occupy the seam block
and whether that block is placed at the right or left join. Hence each panel
has one target list and one permutation to that list. Exterior antisymmetry
contributes exactly the sign of that permutation. If repeated factors occur,
the panel is zero and is outside the nonzero sole-seam case.
\end{proof}

\begin{theorem}[Grade-\(r\) sole-seam identity]
\label{thm:grade-r-seam}
Under these hypotheses,
\[
  (A\mathbin{\lrcorner}B)\wedge(A\mathbin{\lrcorner}C)
  =
  (A\mid S)\,
  \bigl(A\mathbin{\lrcorner}(B\sqcup C)\bigr).
\]
\end{theorem}

\paragraph{Operational intuition (not part of the proof).}
Sort each panel using any fixed total order, turn the two seam blocks toward
the join, and pay the two permutation parities. Stitch the aligned panels,
then apply the grade gate: descend when the accumulated seam is still larger,
emit the overlap scalar when grades balance, and return zero when alternation
forces a repeated residual seam.

\begin{proof}
Write \(A=a_1\wedge\cdots\wedge a_r\) and
\(S=s_1\wedge\cdots\wedge s_r\). For an ordered blade
\[
X=x_1\wedge\cdots\wedge x_m
\]
and a position set \(I=\{i_1<\cdots<i_r\}\subseteq\{1,\ldots,m\}\), put
\[
X_I=x_{i_1}\wedge\cdots\wedge x_{i_r},\qquad
X_{\widehat I}=x_1\wedge\cdots\wedge\widehat{x_{i_1}}\wedge\cdots
\wedge\widehat{x_{i_r}}\wedge\cdots\wedge x_m.
\]
With the contraction convention of Section 1, define
\[
\operatorname{sgn}(I):=(-1)^{\,i_1+\cdots+i_r-r}.
\]
This is the sign of moving the selected factors to the contraction slots,
combined with the reversal forced by the iterated left contractions. Thus
\[
\Delta_A(I;X)
:=\operatorname{sgn}(I)
\det\!\left[(a_\alpha\cdot x_{i_\beta})\right]_
             {\alpha,\beta=1}^{r}
\]
and the \(r\)-fold contraction is the explicit expansion
\[
A\mathbin{\lrcorner}X
 =\sum_{\lvert I\rvert=r}\Delta_A(I;X)\,X_{\widehat I}.
\tag{3.1}
\]
For \(r=1\), this gives the usual sign
\((-1)^{i_1-1}\). For \(I=\{1,\ldots,r\}\), it gives the determinant
pairing with the displayed ordered block, including the sign dictated by the
contraction convention.

We now compare the coefficients of each residual wedge. Write
\[
B=b_1\wedge\cdots\wedge b_p\wedge s_1\wedge\cdots\wedge s_r,\qquad
C=s_1\wedge\cdots\wedge s_r\wedge c_1\wedge\cdots\wedge c_q,
\]
temporarily taking \(\epsilon_B=\epsilon_C=1\). Let
\(P\subseteq\{1,\ldots,p\}\), \(Q\subseteq\{1,\ldots,q\}\), and
\(R\subseteq\{1,\ldots,r\}\) describe a residual wedge
\[
Y_{P,R,Q}:=b_P\wedge s_R\wedge c_Q,
\]
where each subscript retains the original order. A term on the left can
produce this residual only by partitioning the residual seam indices as
\(R=T\mathbin{\dot\cup}(R\setminus T)\): the factors \(s_T\) come from the
first panel and \(s_{R\setminus T}\) from the second. Define
\[
\begin{aligned}
I_{P,T}
  &:=\bigl(\{1,\ldots,p\}\setminus P\bigr)
    \cup\bigl(p+\bigl(\{1,\ldots,r\}\setminus T\bigr)\bigr),\\
J_{T,Q}
  &:=\bigl(\{1,\ldots,r\}\setminus(R\setminus T)\bigr)
    \cup\bigl(r+\bigl(\{1,\ldots,q\}\setminus Q\bigr)\bigr),\\
K_{P,R,Q}
  &:=\bigl(\{1,\ldots,p\}\setminus P\bigr)
    \cup\bigl(p+\bigl(\{1,\ldots,r\}\setminus R\bigr)\bigr)\\
  &\qquad\cup\bigl(p+r+\bigl(\{1,\ldots,q\}\setminus Q\bigr)\bigr).
\end{aligned}
\]
Only choices with all three displayed index sets of cardinality \(r\)
contribute. The seam shuffle has sign
\[
\eta(R,T)
:=(-1)^{\#\{(u,v)\in T\times(R\setminus T):u>v\}},
\]
because \(s_T\wedge s_{R\setminus T}=\eta(R,T)s_R\). Therefore the
coefficient of \(Y_{P,R,Q}\) on the left is
\[
\sum_T\eta(R,T)\,
\Delta_A(I_{P,T};B)\,
\Delta_A(J_{T,Q};C).
\tag{3.2}
\]

The Laplace--Cauchy--Binet minor identity, applied to the two copies of the
ordered seam block \(S\), gives
\[
\sum_T\eta(R,T)\,
\Delta_A(I_{P,T};B)\,
\Delta_A(J_{T,Q};C)
=(A\mid S)\,\Delta_A(K_{P,R,Q};b\wedge S\wedge c).
\tag{3.3}
\]
The sign comparison used in this convolution is the following elementary
index calculation.

\begin{lemma}[Sign in the minor convolution]
Let \([n]=\{1,\ldots,n\}\), write \(t=\lvert T\rvert\), \(k=\lvert R\rvert\),
and put
\[
h(R,T):=\#\{(u,v)\in T\times(R\setminus T):u>v\}.
\]
Let
\[
\operatorname{sgn}_S:=(-1)^{\binom r2}
\]
be the sign of the ordered seam block considered by itself. Then
\[
\operatorname{sgn}(I_{P,T})\operatorname{sgn}(J_{T,Q})\eta(R,T)
=(-1)^{p(k-t-q+\lvert Q\rvert)+h(R,T)}
\operatorname{sgn}_S\operatorname{sgn}(K_{P,R,Q}).
\tag{3.4}
\]
\end{lemma}

\begin{proof}
Write
\[
\begin{aligned}
\sum I_{P,T}
 &=\sum([p]\setminus P)+p(r-t)+\sum([r]\setminus T),\\
\sum J_{T,Q}
 &=\sum([r]\setminus R)+\sum T
   +r(q-\lvert Q\rvert)+\sum([q]\setminus Q),\\
\sum K_{P,R,Q}
 &=\sum([p]\setminus P)+p(r-k)+\sum([r]\setminus R)\\
 &\qquad +(p+r)(q-\lvert Q\rvert)+\sum([q]\setminus Q).
\end{aligned}
\]
Using \(\operatorname{sgn}(I)=(-1)^{\sum I-r}\), subtraction gives
\[
\begin{aligned}
&\bigl(\sum I_{P,T}-r\bigr)
 +\bigl(\sum J_{T,Q}-r\bigr)
 -\bigl(\sum K_{P,R,Q}-r\bigr)
 +h(R,T)\\
&\qquad\equiv
p(k-t-q+\lvert Q\rvert)+\binom r2+h(R,T)
\pmod 2.
\end{aligned}
\]
Since \(\operatorname{sgn}_S=(-1)^{\binom r2}\), this is exactly (3.4).
The first term in the remaining exponent counts the index shifts across the
\(b\)-block and the \(c\)-block; \(h(R,T)\) counts the seam shuffle needed to
restore \(s_T\wedge s_{R\setminus T}\) to \(s_R\). These are precisely the
cofactor and shuffle signs in the Laplace expansion. Hence each \(T\)-term
has the sign on the right-hand side of (3.3), and the ordinary
Laplace--Cauchy--Binet minor identity proves (3.3).
\end{proof}

Thus (3.3) is not relying on an unspecified ``prescribed'' sign: the
selected-index signs, block shifts, and seam shuffle have all been reduced to
the explicit parity in (3.4).

\subsection{The Laplace identity being used}

The determinant result invoked here is the generalized Laplace expansion for
complementary minors. For an \(n\times n\) matrix \(M\), and fixed row and
column sets \(I,J\subseteq\{1,\ldots,n\}\) of size \(k\),
\[
\det M
 =\sum_{\lvert L\rvert=k}
   (-1)^{\sum I+\sum L}
   \det M[I,L]\det M[I^c,L^c],
\tag{3.5}
\]
where the complementary row set is \(I^c\), the complementary column set is
\(L^c\), and all minors use increasing index order. This is the generalized
Laplace theorem, equivalently the Cauchy--Binet identity applied to compound
matrices; see Gantmacher~[1959, Ch.~I, \S5].

For the present rectangular pairing matrix, the exact compound-minor
specialization is obtained by putting
\[
[I]_Y:=\det\!\left[(a_\alpha\cdot y_{i_\beta})\right]_
             {\alpha,\beta=1}^{r}.
\]
With the \(b,S,c\) column blocks and the index sets defined above, (3.5)
becomes
\[
\sum_T\eta(R,T)\,
\operatorname{sgn}(I_{P,T})\operatorname{sgn}(J_{T,Q})
[I_{P,T}]_B[J_{T,Q}]_C
=\operatorname{sgn}_S\operatorname{sgn}(K_{P,R,Q})
[K_{P,R,Q}]_{b\wedge S\wedge c}[S]_S.
\tag{3.6}
\]
The sign lemma verifies the conversion from the increasing-column minors in
(3.5) to the contraction minors \(\Delta_A\). Since
\(A\mid S=\operatorname{sgn}_S[S]_S\), equation (3.6) is exactly (3.3), not
an additional heuristic sign assertion.

\subsection{Recovery of the base case}

Set \(r=1\), \(p=q=1\), \(A=a\), \(S=s\), \(b=b\), and \(c=c\). Then
\[
a\mathbin{\lrcorner}(b\wedge s)=(a\cdot b)s-(a\cdot s)b,
\qquad
a\mathbin{\lrcorner}(s\wedge c)=(a\cdot s)c-(a\cdot c)s.
\]
Their wedge is
\[
\begin{aligned}
&(a\mathbin{\lrcorner}(b\wedge s))\wedge
  (a\mathbin{\lrcorner}(s\wedge c))\\
&=(a\cdot s)\bigl[
 (a\cdot b)s\wedge c-(a\cdot s)b\wedge c
 +(a\cdot c)b\wedge s\bigr]\\
&=(a\cdot s)\bigl(a\mathbin{\lrcorner}(b\wedge s\wedge c)\bigr).
\end{aligned}
\]
Here the term containing \(s\wedge s\) is the repeated-residual term that
vanishes. Since \(\operatorname{sgn}(\{1\})=1\) and
\(\operatorname{sgn}_S=1\), this is exactly the sole-shared-seam identity of
Section 2, including its sign convention.

There is an important qualification about vanishing. If a seam factor
\(s_k\) is omitted from both selected sets, then the residual wedge contains
\(s_k\wedge s_k\) and is zero by alternation. However, a term with
\(I\ne I_S\) need not vanish individually: its omitted seam factors may be
complementary to those omitted in the other panel. Such nonprincipal terms
are the \(T\)-summands in (3.2), and they combine by (3.3). It would therefore
be incorrect to discard every \(I\ne I_S\) term; only repeated-residual
terms vanish before the minor convolution.

Finally, (3.1) applied to \(b\wedge S\wedge c\), followed by (3.3), shows
\[
\begin{aligned}
(A\mathbin{\lrcorner}B)\wedge(A\mathbin{\lrcorner}C)
 &= (A\mid S)\bigl(A\mathbin{\lrcorner}(b\wedge S\wedge c)\bigr)\\
 &= (A\mid S)\bigl(A\mathbin{\lrcorner}(B\sqcup C)\bigr).
\end{aligned}
\]
Restoring the orientation factors multiplies both sides by
\(\epsilon_B\epsilon_C\), which is already part of \(B\sqcup C\).
\end{proof}

\begin{remark}
The grade count explains why the seam must have grade \(r\). If
\(\operatorname{grade}(A)=r\), \(\operatorname{grade}(S)=t\), and
\(\operatorname{grade}(b)=p\), \(\operatorname{grade}(c)=q\), then the two
sides before the coefficient have grades
\[
  p+t-r+q+t-r
  \quad\text{and}\quad
  p+t+q-r,
\]
respectively. Their difference is \(t-r\). Thus a scalar coefficient is
possible exactly when \(t=r\). The vector proposition is the \(r=t=1\)
case; the recursive coefficient uses the grade-\(r\) seam identity.
\end{remark}

\section{Nested Recursion}

\subsection{Orientation and parity}

Before applying the seam identity, each panel is oriented so that the seam is
rightmost in the left panel and leftmost in the right panel:
\[
B=\epsilon_B\,b\wedge S,\qquad
C=\epsilon_C\,S\wedge c,\qquad
\epsilon_B,\epsilon_C\in\{+1,-1\}.
\]
The factors \(\epsilon_B\) and \(\epsilon_C\) record the parity of these
reorderings. They are part of the coefficient, so alignment is mathematical
data rather than a cosmetic normalization step.

The grade-\(r\) identity resolves the local seam step. The remaining nested
construction is an iteration of vector contractions, not an application of
the vector Leibniz rule to an accumulated blade. Write
\[
  \iota_u(X):=u\mathbin{\lrcorner}X .
\]
For an ordered chain of vectors \(d_1,\ldots,d_r\), define
\[
\begin{aligned}
A_{\mathrm{acc}}^{(r)}
  &:=d_1\wedge d_2\wedge\cdots\wedge d_r,\\
\iota_{A_{\mathrm{acc}}^{(r)}}(X)
  &:=\iota_{d_1}\bigl(
      \iota_{d_2}(\cdots\iota_{d_r}(X)\cdots)
    \bigr).
\end{aligned}
\]
The recursive contraction rule
\[
  (u\wedge A)\mathbin{\lrcorner}X
  =
  u\mathbin{\lrcorner}(A\mathbin{\lrcorner}X)
\]
shows that this is exactly the corresponding blade contraction under the
chosen ordering. Each local application of Theorem
\ref{thm:grade-r-seam} is made with the current grade-\(r\) accumulated
contractor and its grade-\(r\) accumulated seam. No blade-level Leibniz rule
is being assumed.

At depth \(k\), two nested panels may be written schematically as
\[
\begin{aligned}
T_L&=\iota_{d_1}\left(
 X^L_1\wedge\iota_{d_2}\left(
 X^L_2\wedge\cdots\iota_{d_k}(C_L)\cdots\right)\right),\\
T_R&=\iota_{d_1}\left(
 X^R_1\wedge\iota_{d_2}\left(
 X^R_2\wedge\cdots\iota_{d_k}(C_R)\cdots\right)\right).
\end{aligned}
\]
When the recursive seam has accumulated to
\(S_{\mathrm{acc}}^{(r)}=s_1\wedge\cdots\wedge s_r\), the coefficient
\[
  A_{\mathrm{acc}}^{(r)}\mid S_{\mathrm{acc}}^{(r)}
\]
is scalar because both blades have grade \(r\). The matching growth of the
contractor and seam is therefore forced by grade, rather than an arbitrary
convention.

At a given level, let
\[
x=
\operatorname{grade}(D_L)
+\operatorname{grade}(D_R)
-\operatorname{grade}(D_{\mathrm{merged}})
-\operatorname{grade}(A_{\mathrm{acc}}).
\]
Under the sole-seam invariant, the three cases are:
\[
\begin{array}{c|l}
x<0 & \text{recurse into the merged down expression},\\
x=0 & \text{apply the grade-\(r\) identity and emit its scalar},\\
x>0 & \text{the expression vanishes because residual seam factors repeat}.
\end{array}
\]

The local grade-\(r\) seam identity is now a theorem. The remaining global
claim concerns repeated application through arbitrary nested trees:

\begin{quote}
\textbf{Inductive invariant.}
Let \(h(\mathcal T)\) be the number of unresolved recursive seam levels in a
pair of nested panel trees. We use the following level-local invariant:
\begin{enumerate}
  \item every down expression at a node is decomposable and homogeneous;
  \item after orientation, the two panels have the form
        \(B=b\wedge S\) and \(C=S\wedge c\), with \(S\) their sole common
        visible factor;
  \item a nested argument is opaque at its parent node and is inspected only
        when recursion reaches that argument's own node; and
  \item all scalar factors and orientation parities produced above a node are
        carried unchanged into its merged child.
\end{enumerate}
The word ``visible'' refers to immediate exterior-product arguments, not to
the labelled leaves occurring inside nested contractions.
\end{quote}

\subsection{Level-local visibility}

The closure argument uses a syntactic, rather than a global support,
property. The panel grammar has three relevant constructors:
\[
E ::= \text{atom}\mid E_1\wedge\cdots\wedge E_n
       \mid (U\mathbin{\lrcorner}D).
\]
At a fixed exterior-product node \(D=E_1\wedge\cdots\wedge E_n\), define
\[
\operatorname{Vis}(D):=\{E_1,\ldots,E_n\}.
\]
The seam matcher compares elements of \(\operatorname{Vis}(D_L)\) and
\(\operatorname{Vis}(D_R)\); it does not compare arbitrary labelled leaves
inside a nested contraction. A proper descendant of \(E_i\) is therefore
opaque at this node. If the whole nested expression \(E_i\) occurs on both
sides, it is one visible seam, and its own descendants are inspected only
when recursion reaches that child node.

\begin{lemma}[Syntactic visibility of seam candidates]
\label{lem:syntactic-visibility}
Let \(q\) be a proper descendant of an immediate factor \(E_i\) at an exterior
node \(D\). If \(q\ne E_i\), then \(q\notin\operatorname{Vis}(D)\) and cannot
be a seam candidate at \(D\). A repeated nested expression \(E_i\), by
contrast, is one candidate at \(D\); its internal factors are candidates only
at the nested node that contains them.
\end{lemma}

\begin{proof}
By definition, \(\operatorname{Vis}(D)\) contains exactly the immediate
arguments of the exterior node. A proper descendant of \(E_i\) is an argument
of a child of \(E_i\), not an argument of \(D\), so it is not in
\(\operatorname{Vis}(D)\). The seam matcher therefore cannot select it.

If the complete expression \(E_i\) occurs on both sides, the two immediate
argument lists contain the same element \(E_i\), so they expose one seam.
The recursive matcher receives \(E_i\) as its own node only after descending
to that node and then uses \(\operatorname{Vis}(E_i)\). No descendant is
lifted into the parent argument list.
\end{proof}

The following lemma supplies the algebraic bridge from visibility to the
permitted expansion.

\begin{lemma}[Semantic sibling separation]
\label{lem:semantic-sibling-separation}
Let \(D=E_1\wedge\cdots\wedge E_n\), where each \(E_i\) is either a one-blade
or a nested contraction node, and let \(A\) be a blade of grade \(r\). In the
full multilinear expansion of \(A\mathbin{\lrcorner}D\), expanding every
nested node down to leaves, each monomial has the form
\[
\pm\,\Delta\cdot(F_1\wedge\cdots\wedge F_n),
\]
where \(F_i\) is a fully expanded residual built only from the leaves in the
subtree \(E_i\), and \(\Delta\) is a product of pairings. No pairing term uses
down-side leaves from two distinct sibling subtrees \(E_i\) and \(E_j\).
\end{lemma}

\begin{proof}
For a flat exterior node, write \(A=a_1\wedge\cdots\wedge a_r\). Repeated
application of the graded Leibniz rule partitions the ordered contractor
factors among the siblings:
\[
A\mathbin{\lrcorner}(E_1\wedge\cdots\wedge E_n)
=\sum_{I_1\mathbin{\dot\cup}\cdots\mathbin{\dot\cup}I_n=[r]}
\varepsilon(I_1,\ldots,I_n)
\bigwedge_{i=1}^{n}(A_{I_i}\mathbin{\lrcorner}E_i),
\]
where \(A_{I_i}\) retains the original order. To make the sign explicit, let
\(\iota(j)=i\) when \(j\in I_i\), let
\(g_i=\operatorname{grade}(E_i)\), and put
\[
\varepsilon(I_1,\ldots,I_n)
:=(-1)^{\,
 \sum_{j=1}^{r}\sum_{k<\iota(j)}g_k
 +\#\{(j,\ell):j<\ell,\ \iota(j)>\iota(\ell)\}}.
\]
The first term is the sum of the grades crossed by the \(j\)-th contraction;
the second is the inversion parity created when an earlier contraction lowers
the grade crossed by a later one. For a one-blade \(E_i\), the factor is
zero unless the assigned grade is admissible, and its nonzero coefficient is
the corresponding dot pairing. For a nested \(E_i\), the same expression is
expanded recursively inside that factor. Induction on the total number of
nested contraction nodes therefore expands each
\(A_{I_i}\mathbin{\lrcorner}E_i\) independently. All pairings created there
have both their down-side leaves inside \(E_i\); wedging the residuals
introduces only signs and possible repeated-factor zeros, never a new pairing
across siblings. This proves the claim.
\end{proof}

The syntactic visibility lemma is therefore a statement about what the
matcher can select, while semantic sibling separation supplies the algebraic
bridge. Together they justify ignoring a buried coincidence without defining
it away: Theorem~\ref{thm:nested-induction} proves that the permitted stitches
still equal full direct left-contraction expansion. In particular, this covers
the hidden-coincidence configuration
\[
D_L=a_2\wedge a_3\wedge
 \bigl(a_4\mathbin{\lrcorner}(b_1\wedge b_2)\bigr),
\qquad D_R=a_3\wedge b_1,
\]
which is also checked by the corresponding implementation regression test.

\begin{lemma}[One-stitch closure]
\label{lem:one-stitch-closure}
Let \(T_L,T_R\) be well-formed nested contraction trees of height \(h\)
satisfying the level-local sole-seam invariant and sharing a contraction
lineage. For the stitch defined in the Algorithm section, grammar preservation
follows directly from its definition: it concatenates immediate exterior
arguments and recursively visits a child only at that child's own node. If
\(T_M\) is the result of one valid root stitch, then \(T_M\) satisfies the same
invariant at every remaining level.
\end{lemma}

\begin{proof}
Induct on \(h\). For \(h=0\), there are no remaining levels.

For \(h>0\), orient the root down expressions as
\[
D_L=b\wedge S,\qquad D_R=S\wedge c.
\]
The stitch constructs
\[
D_M=b\wedge S\wedge c,
\]
with \(S\) counted once. At this node, the only possible seam candidates
come from the immediate argument lists. Lemma~\ref{lem:syntactic-visibility}
excludes every factor buried inside a nested argument. Thus a factor of \(c\)
that coincides with a buried factor of \(T_L\) cannot create a second seam at
the root. If it coincides with a complete immediate nested argument, that
complete argument is visible and the original sole-seam hypothesis already
excludes a second candidate.

By Lemma~\ref{lem:semantic-sibling-separation} and the definition of the
stitch, \(c\) is not paired with a buried leaf of \(b\)'s nested child, nor is
a child factor lifted into the parent list. Hence every remaining comparison is
either the unique corresponding comparison inside \(S\), or a comparison at a
separate child node whose immediate arguments were already governed by the
sole-seam invariant. The former has strictly smaller height and is preserved
by the induction hypothesis; the latter cannot acquire a new cross-level
candidate by Lemma~\ref{lem:syntactic-visibility}.

Therefore the merged tree has exactly one active seam at every remaining
lineage level. Reordering signs and scalar overlap coefficients change
coefficients only, not the syntax or visibility relation.
\end{proof}

The earlier hereditary-freshness condition is sufficient but stronger than
necessary: labelled leaves may coincide across levels, as long as nested
constructions remain opaque to seam matching at their parent node. The lemma
therefore applies unconditionally to the stated well-formed panel grammar.

\begin{theorem}[Nested sole-seam recursion under the invariant]
\label{thm:nested-induction}
For a fixed valid merge schedule, any pair of well-formed nested contraction
trees satisfying the level-local sole-seam invariant is reduced by repeated
grade-gated stitching to the same expression as direct left-contraction
expansion. The result includes all overlap scalars and exterior-algebra parity
signs.
\end{theorem}

\begin{proof}
We induct on \(h(\mathcal T)\). If \(h(\mathcal T)=0\), the current seam is
grade-matched with the accumulated contractor, so \(x=0\). The
grade-\(r\) sole-seam identity, Theorem~\ref{thm:grade-r-seam}, gives exactly
the direct contraction expansion, including the orientation factor in
\(B\sqcup C\).

Now assume the claim for all trees of height at most \(h\), and consider a
tree of height \(h+1\). Orient the root panels and expand each contraction by
the graded Leibniz rule. If \(x>0\), Lemma~\ref{lem:semantic-sibling-separation}
says that contractor factors assigned to one sibling cannot remove a seam
factor through a pairing with a buried leaf in another sibling. Since the
visible seam has grade larger than the available contractor grade, every
summand retains a common seam factor in both residual wedges. Each such
summand is zero by alternation, and the grade-gated recursion returns zero as
well.

If \(x=0\), the current contractor and seam have equal grade. The root is
therefore exactly the grade-\(r\) situation of
Theorem~\ref{thm:grade-r-seam}; replacing the two panels by their merged
panel preserves the direct expansion and extracts the scalar
\(A_{\mathrm{acc}}\mid S_{\mathrm{acc}}\).

It remains to consider \(x<0\). The graded Leibniz expansion groups the
nonzero terms by the unique residual visible seam. By the one-stitch closure
lemma, this group is represented by a merged down-expression of height at
most \(h\), with exactly the scalar and parity factors recorded at the root.
The induction hypothesis identifies the recursive reduction of that merged
expression with its direct left-contraction expansion. Multiplying by the
unchanged root factors gives the direct expansion of the original tree. This
proves the claim for height \(h+1\), and hence for every finite height.
\end{proof}

The theorem is deliberately stated for a fixed valid merge schedule. The
closure lemma establishes its level-local structural hypothesis for the
well-formed panel grammar. The theorem does not claim that different valid
schedules are confluent. Computational checks of nesting depths up to six,
including higher-grade decomposable blades, support the invariant but do not
replace its proof.

\subsection{The overlap coefficient via Capelli}

At a balanced grade gate, the coefficient is the scalar pairing of the
accumulated contractor and accumulated seam:
\[
  A_{\mathrm{acc}}\mid S_{\mathrm{acc}}
  =
  \det\bigl[(d_i\cdot s_j)\bigr]_{i,j=1}^{r},
\]
with the determinant sign interpreted in the chosen contraction convention.
This determinant is the finite Capelli contraction of the accumulated
contractor with the grade-matched seam. Reordering parities are multiplied
into the coefficient separately.

\section{Algorithm (Stitching)}

The rewrite is described abstractly:

\begin{lstlisting}
canonicalize(expression):
    while a mergeable pair exists:
        choose two panels with a common contraction lineage
        orient the panels around one shared seam
        record the two reordering parities
        split their down blades into unique parts and one seam
        accumulate the contractor and seam blades
        compute the grade-balance value x
        if x > 0:
            replace the expression by zero under the sole-seam invariant
        else if x < 0:
            recursively canonicalize the merged down expression
        else:
            multiply by the scalar contractor-seam pairing and both parities
            replace the pair by the merged contraction
    return expression
\end{lstlisting}

\section{CGA encoding}

The applications use the standard conformal model rather than introducing a
new encoding. For \(x\in\mathbb R^n\), its conformal point is the null vector
\[
\dot{x}=x+e_0+\frac{x^2}{2}e_\infty,
\qquad
\dot{x}^{\,2}=0,
\]
with the usual normalization \(e_0\cdot e_\infty=-1\). A circle or sphere is
represented by the blade obtained from its defining points or carrier
objects. We write \(\vee\) for the standard dual/reduced meet. Incidence is
expressed by
\[
A\wedge\dot y=0.
\]
For a sphere with Euclidean centre \(c\) and radius \(\rho\), the dual
representation is
\[
\dot c-\frac{\rho^2}{2}e_\infty.
\]
These conventions, including lines, circles, spheres, duality, and incidence,
are standard; see Li~\cite{li2026}, Doran and Lasenby~\cite{doran2003}, and
Dorst, Fontijne, and Mann~\cite{dorst2007}. The present paper only applies
stitching to the resulting exterior-product and contraction expressions.

\subsection{Radius-augmented Lie-sphere chart}

The paper uses two representation models. The preceding point-only conformal
model is used for point, line, circle, and sphere incidence. For oriented
circles in the Monge application, we additionally use the radius-augmented
Lie-sphere chart
\(\mathbb R^{3,2}\), following the standard Lie-sphere framework
\cite{cecil1992}. Take basis
\(e_0,e_x,e_y,e_r,e_\infty\) with
\[
 e_0\cdot e_\infty=-1,\qquad
 e_x^2=e_y^2=1,\qquad e_r^2=-1,
\]
and all other products zero. An oriented circle with centre \((x,y)\) and
signed radius \(r\) is
\[
 \mathcal C(x,y,r)
 :=e_0+x e_x+y e_y+r e_r
   +\frac{x^2+y^2-r^2}{2}e_\infty .
\]
It is null, since
\[
 \mathcal C\cdot\mathcal C
 =x^2+y^2-r^2-(x^2+y^2-r^2)=0.
\]
The \(e_r\)-coordinate carries the signed radius. This chart is used only
where the radius coordinate is geometrically relevant; the two models share
the same exterior and contraction operations.

\section{Applications}

The step counts below count successful seam-resolution steps, not elementary
geometric constructions, scalar products, or final incidence tests.

\subsection{A parameterized two-stitch incidence family}

The Monge, perpendicular-bisector, and angle-bisector constructions have the
same algebraic shape when their chart encodings satisfy the same panel
hypotheses. Let \(X_1,X_2,X_3\) be oriented carrier blades and let \(K\) be a
fixed contractor. Put
\[
Q_{ij}:=K\mathbin{\lrcorner}(X_i\wedge X_j).
\]
Assume that the three panels are oriented so that \(Q_{12}\) and \(Q_{23}\)
have the sole seam \(X_2\), that \(Q_{12}\wedge Q_{23}\) exposes the
two-factor seam \(X_3\wedge X_1\) against \(Q_{31}\), and that
\(\operatorname{grade}(K)=\operatorname{grade}(X_2)=1\). The same two-stitch
proof then applies to every
geometric interpretation of the carriers.

\begin{proposition}[Parameterized pairwise-difference incidence]
\label{prop:pairwise-difference}
Under the preceding panel hypotheses, the three pairwise objects \(Q_{12}\),
\(Q_{23}\), and \(Q_{31}\) satisfy
\[
Q_{12}\wedge Q_{23}\wedge Q_{31}=0.
\]
Consequently, whenever the resulting blades represent the relevant affine
objects in a common chart, their three pairwise constructions have the
corresponding common incidence or collinearity relation.
\end{proposition}

\begin{proof}
The first seam is grade-matched, so the grade-\(r\) identity gives
\[
Q_{12}\wedge Q_{23}
=(K\mid X_2)\,
K\mathbin{\lrcorner}(X_1\wedge X_2\wedge X_3),
\]
up to the canonical orientation sign already fixed by the panel order.
Against \(Q_{31}\), the exposed seam is \(X_3\wedge X_1\), of grade two,
whereas \(K\) has grade one. The grade gate is therefore in its overflow
branch: every full-expansion term retains a repeated seam factor and vanishes
by alternation. Hence
\[
Q_{12}\wedge Q_{23}\wedge Q_{31}=0.
\]
The proof depends only on the contractor grade, seam pattern, and orientation;
the geometric meaning of \(X_i\) is the parameter.
\end{proof}

The intended instantiations are summarized below. The first row is realized in
the Lie-sphere chart of this paper. The other two rows are algebraic targets
for a chart-specific encoding, not claims that the point-only CGA chart
automatically supplies those panels.
\[
\begin{array}{c|c|c|c|c}
\text{construction} & K & \text{carriers} & \text{pairwise result}
& \text{status here}\\
\hline
\text{Monge} & e_r & \text{oriented circles}
& \text{homothety centre} & \text{proved}\\
\text{perpendicular bisectors} & e_\infty & \text{points}
& \text{perpendicular bisector} & \text{encoding required}\\
\text{angle bisectors} & e_r & \text{lines}
& \text{angle bisector} & \text{encoding required}
\end{array}
\]
Thus the proposition expresses the common two-stitch architecture, while the
carrier chart must still be checked for each construction. In particular, the
often-suggested line representative
\[
L=x e_x+y e_y+e_r+d\,e_\infty,\qquad x^2+y^2=1,
\]
is null in the radius-augmented metric, but it does not make the angle-bisector
row valid by itself. For
\[
L_1=e_x+e_r+2e_\infty,\qquad
L_2=e_y+e_r-e_\infty,
\]
the direct metric calculation gives
\[
L_1^2=L_2^2=0,\qquad
e_r\mathbin{\lrcorner}(L_1\wedge L_2)=L_1-L_2,
\qquad
\bigl(L_1-L_2\bigr)^2=2\ne0.
\]
The contracted blade is therefore not null and cannot be identified with a
null line in this naive chart. This numerical check prevents the angle- and
perpendicular-bisector rows from being overstated; a different
representation, or a separate metric lemma, is required before they become
fully verified CGA applications.

\paragraph{Verified instantiation: Monge.}

We now use the Lie-sphere chart introduced above. The sign of \(r\) records
orientation; reversing it exchanges the internal and external homothety
interpretations.

\begin{lemma}[Radius annihilator]
\label{lem:radius-annihilator}
For two oriented circles \(\mathcal C_i,\mathcal C_j\), put
\[
 P_{ij}:=e_r\mathbin{\lrcorner}
       (\mathcal C_i\wedge\mathcal C_j).
 \]
Then
\[
 P_{ij}
 = (e_r\cdot\mathcal C_i)\mathcal C_j
   -(e_r\cdot\mathcal C_j)\mathcal C_i
 =r_j\mathcal C_i-r_i\mathcal C_j.
\]
Its \(e_r\)-coefficient is zero, and its \(e_0\)-coefficient is
\(r_j-r_i\). If \(r_i\ne r_j\), its normalized Euclidean centre is
\[
 \frac{r_j c_i-r_i c_j}{r_j-r_i},
\]
the homothety centre of the two circles. The vector \(P_{ij}\) need not
itself be null; the relevant homogeneous point coordinates are its
\((e_0,e_x,e_y)\)-components.
\end{lemma}

\begin{proof}
The first equality is the vector Leibniz rule, and
\(e_r\cdot\mathcal C_i=-r_i\). The \(e_r\)-coefficient of
\(r_j\mathcal C_i-r_i\mathcal C_j\) is
\(r_jr_i-r_ir_j=0\). Dividing its \(e_0\)-normalized spatial part by
\(r_j-r_i\) gives the displayed centre.
\end{proof}

\begin{theorem}[Monge]
\label{thm:monge}
For three oriented circles in general position, the three homothety centres
\(P_{12},P_{23},P_{31}\) are collinear.
\end{theorem}

\begin{proof}
First stitch the panels \(P_{12}\) and \(P_{23}\). Their down blades have the
sole visible seam \(\mathcal C_2\), already oriented as
\[
\mathcal C_1\wedge\mathcal C_2,\qquad
\mathcal C_2\wedge\mathcal C_3.
\]
The contractor \(e_r\) and seam \(\mathcal C_2\) both have grade one, so the
balanced grade-\(r\) identity gives
\[
\begin{aligned}
P_{12}\wedge P_{23}
 &= (e_r\cdot\mathcal C_2)\,
    e_r\mathbin{\lrcorner}
    (\mathcal C_1\wedge\mathcal C_2\wedge\mathcal C_3)\\
 &=-r_2\,
    e_r\mathbin{\lrcorner}
    (\mathcal C_1\wedge\mathcal C_2\wedge\mathcal C_3).
\end{aligned}
\]
The resulting panel and \(P_{31}\) now expose the two-factor seam
\[
 S=\mathcal C_3\wedge\mathcal C_1,
\]
while the accumulated contractor remains \(e_r\). Hence
\(\operatorname{grade}(S)-\operatorname{grade}(e_r)=2-1>0\).
The grade gate is therefore in its vanishing branch: every term in the full
expansion retains a repeated residual seam factor, and alternation makes it
zero. Consequently,
\[
 P_{12}\wedge P_{23}\wedge P_{31}=0.
\]
Projecting this linear dependence to the \((e_0,e_x,e_y)\)-coordinates gives
the corresponding dependence of the three homogeneous Euclidean point
coordinates. Since the \(e_0\)-coefficients are nonzero in general position,
their normalized centres are collinear.
\end{proof}

The signed-radius representation covers the orientation cases uniformly.
Circles with radii of the same sign yield the external homothety centre,
whereas opposite signs yield the internal one. The proof does not inspect
these signs: it consists of one balanced stitch followed by one grade
overflow. Thus the same two-step reduction proves the classical external
case and all mixed internal/external cases realizable by a consistent signed
orientation. The all-internal choice is not generally collinear and is not
reachable by any consistent choice of signed radii.
For example, with centres \((0,0),(4,0),(1,3)\) and signed radii
\((-1,-2,-3)\), the three centres are
\((-4,0),(10,-6),(-\tfrac12,-\tfrac32)\), whose affine collinearity
determinant is zero. In contrast, forcing the three pairwise internal centres
for the corresponding positive radii gives determinant \(12/5\), confirming
that the all-internal choice is not a theorem covered by signed orientation.

The reduction uses two successful seam-resolution steps: one balanced merge
emitting the scalar \(e_r\cdot\mathcal C_2=-r_2\), followed by one grade
overflow. The largest intermediate is the grade-three circle wedge
\(\mathcal C_1\wedge\mathcal C_2\wedge\mathcal C_3\). The assumptions exclude
equal-radius centres at infinity and the usual linearly dependent
degeneracies.

\subsection{Concyclic four-point boundary case}

The concyclic four-point case uses a different termination mechanism and is
not an instance of Proposition~\ref{prop:pairwise-difference}. For null point
representatives \(\dot p_1,\ldots,\dot p_4\), concyclicity is the incidence
condition
\[
\dot p_1\wedge\dot p_2\wedge\dot p_3\wedge\dot p_4=0.
\]
Here the computation terminates because the four point representatives are
linearly dependent, not because an accumulated seam has exceeded the
contractor grade. This distinction matters: the example shows that the
stitching calculus uses geometric hypotheses such as concyclicity in addition
to grade balance. We therefore record it as a boundary case rather than claim
a two-stitch contraction proof for it.

\section{Related work}

Li's 2026 treatment of null geometric algebra studies null monomials,
centred null binomials, and their geometric interpretation. In particular,
equation (5.1) gives a factorization rule for products of anisotropic null
monomials with a common null endpoint, reducing the product to versor factors
and endpoint factors~\cite{li2026}. That is a null-monomial/geometric-product
normal-form statement. The identity proved here is different: it is the
local exterior-contraction rewrite
\((A\mathbin{\lrcorner}B)\wedge(A\mathbin{\lrcorner}C)\), with a
grade-dependent overlap scalar and explicit seam parity.

Li, Shao, Huang, and Liu's 2014 work on reduction among bracket polynomials
develops an invariant division and an admissible order leading to an
invariant Gröbner-basis theory~\cite{lishao2014}. We checked the available
statement and proceedings description for an overlap with the present
construction. Its invariant division reduces bracket polynomials modulo an
algebraic ordering; it does not state the local contraction-panel rewrite or
the sole-shared-seam identity proved here. The two approaches are therefore
complementary: invariant division addresses global polynomial reduction,
whereas \(\sandhi\) addresses a local geometric contraction pattern.

The 2008 monograph \emph{Invariant Algebras and Geometric Reasoning} contains
the relevant null-bracket-algebra and Euclidean-geometry chapters, alongside
Grassmann--Cayley, conformal, and Clifford-algebra material
\cite{li2008}. Our check was limited to the available table of contents,
bibliographic records, searchable previews, and descriptions of the
``Null Bracket Algebra and Euclidean Geometry'' material; full access to the
book was not available. Those sources establish the surrounding bracket
algebra context, but they were insufficient to verify every chapter-level
identity, so no claim of exhaustive literature exclusion is made.

\section{Limitations}

The grade-\(r\) seam identity is proved for decomposable homogeneous blades
with the stated orientation and sole-seam hypotheses. The fixed-schedule
induction theorem is proved for the well-formed, level-local panel grammar;
schedule-independent normalization remains open. Since panel orientation and
its parity are canonical once the atom order is fixed, the remaining source of
schedule dependence is pair selection: which eligible seam is stitched first.
Cases through nesting depth six, including higher-grade decomposable blades,
have been verified computationally. The paper does not claim a normal-form
theorem for arbitrary mixed-grade or nondecomposable multivectors.

\begin{thebibliography}{9}

\bibitem{li2026}
H.~Li,
\newblock ``Null geometric algebra with geometric interpretation,''
\newblock \emph{Philosophical Transactions of the Royal Society A},
384:20250096, 2026.
\newblock \url{https://doi.org/10.1098/rsta.2025.0096}.

\bibitem{cecil1992}
T.~E.~Cecil,
\newblock \emph{Lie Sphere Geometry: With Applications to Submanifolds},
\newblock Springer-Verlag, 1992.

\bibitem{lishao2014}
H.~Li, C.~Shao, L.~Huang, and Y.~Liu,
\newblock ``Reduction among bracket polynomials,''
\newblock in \emph{Proceedings of ISSAC 2014}, pp.~304--311, 2014.
\newblock \url{https://doi.org/10.1145/2608628.2608630}.

\bibitem{li2008}
H.~Li,
\newblock \emph{Invariant Algebras and Geometric Reasoning},
\newblock World Scientific, 2008.
\newblock \url{https://doi.org/10.1142/6514}.

\bibitem{doran2003}
C.~Doran and A.~Lasenby,
\newblock \emph{Geometric Algebra for Physicists},
\newblock Cambridge University Press, 2003.

\bibitem{dorst2007}
L.~Dorst, D.~Fontijne, and S.~Mann,
\newblock \emph{Geometric Algebra for Computer Science},
\newblock Morgan Kaufmann, 2007.

\bibitem{gantmacher1959}
F.~R.~Gantmacher,
\newblock \emph{The Theory of Matrices}, Vol.~1,
\newblock Chelsea Publishing Company, 1959, Chapter~I, \S5.

\end{thebibliography}

\end{document}
